\chapter{Experimental Methodology}\label{ch:m-methodology}

This chapter defines the experiment behind the comparison of host-driven and device-initiated communication on a GPU-accelerated cluster. It sets out the measurement design and its parameters, the platform and the configured stacks that are actually compared, the timing and statistical protocol, the communication patterns and the application validation built on them, and the boundaries of what the retained evidence supports. The recurring concern is attribution: the evaluated paths differ in compute interface, memory handling, completion semantics, and transport as well as in where communication is invoked, so the design states throughout which of those differences a given comparison holds fixed and which it does not.

\section{Research Planning}\label{sec:research-planning-levels}

The evaluation employs a bottom-up approach, progressing from isolated microbenchmarks to a full application context. This progression is structured in three levels, and the campaign was executed in that order, because each level uses the preceding one to choose what is worth measuring:

\begin{enumerate}
    \item \textbf{Isolated communication microbenchmarks:} \texttt{pingpong}, \texttt{halo\_1d}, \texttt{allreduce}, and \texttt{alltoall} exercise one operation at a time, with no computation to overlap or distort the transfer. They establish the message sizes, topologies, and runtime settings at which a given backend is competitive.
    \item \textbf{Application microbenchmarks:} \texttt{moe} and \texttt{cg\_step} retain the same harness, timing contract, and validation policy, but replace a single operation with a workload's communication sequence. \texttt{moe} adds routing-dependent, variable-count redistribution; \texttt{cg\_step} combines a halo exchange, a local stencil, and two scalar reductions. They test whether the isolated observations survive composition, either with an irregular payload distribution or with local computation.
    \item \textbf{Full \texttt{aCG} solver:} Tests cross-level agreement by introducing sparse partitioning, realistic application dependencies, and convergence criteria.
\end{enumerate}

Of the two application microbenchmarks, \texttt{cg\_step} is the one carried into the cross-level comparison of Section~\ref{sec:acg-method}, because the solver evaluated in this thesis is a CG solver and its reductions can be matched by datatype and payload. \texttt{moe} has no matched application in this campaign; it is measured to test how strongly the ranking obtained from a uniform, fixed-count exchange depends on the payload distribution assumed by that exchange.

Raising the application context in steps is what makes the levels diagnostic: each one adds a single class of application property, so a ranking that changes between two levels identifies which property caused it.

The unit of comparison is a configured runtime stack, not a library. Each stack fixes a compute interface, a communication path, and a completion contract together, and the evaluated device-initiated paths differ from their host-driven counterparts in all three, not only in where a transfer is invoked. The design therefore aims to locate performance regimes rather than to attribute differences to invocation alone; Section~\ref{sec:eval-environment} gives the seven stacks and the differences each comparison leaves uncontrolled.

\subsection{Experiment Parameters}\label{sec:execution-procedure}

One benchmark measurement is defined by the parameters in Table~\ref{tab:m-experiment-parameters}.

\begin{table}[htbp]
\centering
\caption{Parameters defining one benchmark measurement, and the values used in this campaign.}
\label{tab:m-experiment-parameters}
\footnotesize
\begin{tabularx}{\textwidth}{lYl}
\toprule
Parameter & Role in the measurement & Campaign value\\
\midrule
Pattern & Communication operation or sequence under test & Six (Table~\ref{tab:m-benchmark-patterns})\\
Stack & Compute interface and communication path & Seven (Table~\ref{tab:m-evaluated-stacks})\\
Topology & Node and GPU count; one rank drives one GPU & Seven, \texttt{1n1g}--\texttt{8n4g}\\
Payload & Message size or application problem size & Per pattern (Table~\ref{tab:m-benchmark-patterns})\\
\midrule
Warm-up iterations & Untimed executions preceding measurement & 20; 10 for \texttt{cg\_step}\\
Timed iterations & Executions averaged into one sample & 100; 50 for \texttt{cg\_step}\\
\midrule
Trials per allocation & Process launches sharing one topology & 3\\
Independent allocations & Scheduler jobs with separate node assignments & 5\\
\bottomrule
\end{tabularx}
\end{table}

A \emph{cell} is one $(\text{pattern}, \text{stack}, \text{topology})$ combination, swept over the payload axis. Topology is given as $N$ nodes by $G$ GPUs rather than as a rank count, because equal rank counts can use different interconnects; \texttt{1n4g} and \texttt{2n1g} are both measured for that reason. Each pattern has a fixed collective policy, such as whether UCC is enabled, described in Section~\ref{sec:runtime-method}.

\paragraph{Timing.} Each trial first runs untimed warm-up iterations, which absorb one-time costs such as buffer registration, connection establishment, and algorithm selection, and then its timed iterations. A timed iteration covers the complete operation from submission to observed completion; setup and correctness checks are excluded (Section~\ref{sec:completion-method}). The timed iterations of a trial form one sample series and are not treated as independent repetitions.

\paragraph{Independence.} Every cell was run as five independent repetitions. Each repetition is a separate scheduler job that obtains its own allocation, so repetitions share no process state and may differ in nodes and network routes. Within one allocation the benchmark is launched three times. Only allocations are treated as independent samples, and each contributes equally to the reported statistics (Section~\ref{sec:samples-method}).

\paragraph{Exceptions to the defaults.} The retained all-to-all and \texttt{cg\_step} sweeps contain three allocations instead of five, each still with three trials; their counts are reported as measured and not reweighted. \texttt{cg\_step} also uses 50 timed and 10 warm-up iterations instead of 100 and 20.

\subsection{Benchmark Tooling}\label{sec:harness}

The campaign was run with the command-line tooling of the benchmark suite contributed by this thesis (Chapter~\ref{ch}), available at \url{https://github.com/y3rbiadit0/gpu-comm-benchmark}. It covers the workflow from an empty checkout to recorded results, so that every stack is built, launched, and logged in the same way:

\begin{itemize}
    \item \textbf{Initialization.} One command bootstraps the source-built dependencies, including OSHMPI and the two oneCCL backends, in dependency order, and builds all seven stacks from shared build presets. Libraries provided by the site, such as HPC-X MPI, are taken from the system modules.
    \item \textbf{Execution.} Experiments can be submitted for a single cell, for one pattern, or for the complete declared matrix of patterns, stacks, and topologies. Iteration counts, trials, and collective policy are defined once per pattern and apply to every stack, and the resolved configuration of a cell can be inspected before submission.
    \item \textbf{Library versions.} NCCL and NVSHMEM releases are selected with a single variable that gives each version its own installation, build, and results directory, so measurements of different versions cannot be mixed. Source-built dependencies are pinned to a revision and rebuilt when changed.
    \item \textbf{Recording.} Each trial log records the job, nodes, stack, launcher, and communication-relevant environment. Output is kept as a result only if the trial completes successfully within its timeout; failed runs are retained separately for inspection.
    \item \textbf{Portability.} Machine-specific modules and transport settings are kept apart from the experiment definitions, so the same experiments can target another cluster by adding a cluster definition.
\end{itemize}

\subsection{Research Questions and Evidence}

Table~\ref{tab:m-artifact-evidence} connects the research questions to the measurements and implementation evidence. For RQ2, programming requirements are assessed through allocation, remote addressing, dependency handling, completion, and validation of the adapters. This is a qualitative assessment of concrete integration work, rather than a developer-productivity experiment.

\begin{table}[htbp]
\centering
\caption{Evidence used to answer the research questions.}
\label{tab:m-artifact-evidence}
\footnotesize
\begin{tabularx}{\textwidth}{lYY}
\toprule
Question & Evidence & What the comparison establishes\\
\midrule
RQ1 & Primitive and composite measurements; persistent NVSHMEM kernels; monolithic aCG & Favorable regimes of the evaluated execution paths\\
RQ2 & Direct SHMEM measurements; grouped NCCL and OSHMPI adapters; NVSHMEM runtime validation & Complete-path costs and required semantic adaptations\\
RQ3 & Matched reduction costs, solver scaling, phase measurements, and CUDA--SYCL report & Predictive accuracy, limits of halo representation, and implementation performance retention\\
\bottomrule
\end{tabularx}
\end{table}

The analysis distinguishes the archived performance campaign from later implementation changes. In particular, the measured OSHMPI \texttt{cg\_step} stages reductions to host memory without copying the results back to the GPU, while the measured aCG OSHMPI communicator uses MPI scalar reductions. Current source changes implement a complete staged round trip in both programs, but have no validated Leonardo performance results. They are candidates for a matched follow-up, not replacements for the archived measurements.

\section{Platform and Evaluated Configurations}\label{sec:eval-environment}

All measurements use Leonardo's Booster partition at CINECA: four A100 GPUs per node, direct intra-node NVLink connections, and an InfiniBand HDR network, as described in Section~\ref{sec:leonardo}. One process controls one GPU. A topology label \texttt{$N$n$G$g} denotes $N$ nodes with $G$ ranks and GPUs per node. The campaign spans one to 32 GPUs.

\subsection{Runtime Stacks and Software Versions}

The benchmark compares seven source implementations or backend builds (Table~\ref{tab:m-evaluated-stacks}). The two oneCCL paths compile the same SYCL benchmark source against different backends; Section~\ref{sec:oneccl-method} describes what their adapters add to the measured path.

\begin{table}[htbp]
\centering
\caption{Benchmark runtime identifiers and their communication paths.}
\label{tab:m-evaluated-stacks}
\footnotesize
\begin{tabularx}{\textwidth}{lYY}
\toprule
Identifier & Compute interface & Communication path\\
\midrule
\texttt{cuda\_mpi} & CUDA & GPU-aware HPC-X MPI\\
\texttt{sycl\_mpi} & SYCL & GPU-aware HPC-X MPI over USM\\
\texttt{cuda\_nccl} & CUDA & NCCL\\
\texttt{cuda\_nvshmem} & CUDA & NVSHMEM\\
\texttt{oshmpi} & CUDA & OSHMPI over HPC-X MPI\\
\texttt{sycl\_oneccl} & SYCL & oneCCL/NCCL\\
\texttt{sycl\_oneccl\_oshmpi} & SYCL & oneCCL/OSHMPI\\
\bottomrule
\end{tabularx}
\end{table}

Table~\ref{tab:m-pattern-coverage} records measured coverage rather than API capability. The seven stack identifiers describe the union of the campaign; no single pattern contains all paths. In particular, oneCCL/OSHMPI is measured only for the standalone collectives, while the persistent device-invoked NVSHMEM implementation is specific to the neighbor patterns and the monolithic solver. The MoE column is complete for the six stacks that declare it: oneCCL/OSHMPI is not declared for MoE, because the adapter's point-to-point path would stage every variable-count block through symmetric memory.

\begin{table}[htbp]
\centering
\caption{Measured backend coverage in the reconciled benchmark archive. ``No'' denotes an absent campaign cell, not an unsupported operation in the underlying library.}
\label{tab:m-pattern-coverage}
\scriptsize
\begin{tabular}{lcccccc}
\toprule
Stack & Ping-pong & Halo & Allreduce & All-to-all & \texttt{moe} & \texttt{cg\_step}\\
\midrule
\texttt{cuda\_mpi} & Yes & Yes & Yes & Yes & Yes & Yes\\
\texttt{sycl\_mpi} & Yes & Yes & Yes & Yes & Yes & Yes\\
\texttt{cuda\_nccl} & Yes & Yes & Yes & Yes & Yes & Yes\\
\texttt{cuda\_nvshmem} & Yes & Yes & Yes & Yes & Yes & Yes\\
\texttt{oshmpi} & Yes & Yes & Yes & Yes & Yes & Yes\\
\texttt{sycl\_oneccl} & Yes & Yes & Yes & Yes & Yes & Yes\\
\texttt{sycl\_oneccl\_oshmpi} & No & No & Yes & Yes & No & No\\
\bottomrule
\end{tabular}
\end{table}

\begin{table}[htbp]
\centering
\caption{Software versions and roles in the main campaign. Source-built adapters are identified by repository and branch in the artifact records; exact measured revisions are not available for every archived run.}
\label{tab:m-eval-versions}
\footnotesize
\begin{tabularx}{\textwidth}{lYY}
\toprule
Component & Version or configuration & Role\\
\midrule
Native toolchain & NVHPC 24.5, CUDA 12.4 & CUDA benchmarks and aCG\\
SYCL toolchain & DPC++ 6.3, GCC 12.2, CUDA 12.2 & SYCL benchmarks and solver\\
HPC-X & 2.19; Open MPI over UCX & Direct MPI and OSHMPI substrate\\
UCC & 1.4.0 & Optional MPI collective component\\
NCCL & 2.18.5 & Direct and oneCCL collectives\\
NVSHMEM & 2.11.0, \texttt{ibrc} & Host- and device-invoked SHMEM\\
Intel MPI & Bundled with the oneCCL/NCCL build & Launch and bootstrap; NCCL carries payload\\
OSHMPI/oneCCL & Source-built extensions & Symmetric allocation and backend adaptation\\
\bottomrule
\end{tabularx}
\end{table}

MPI and OSHMPI share an MPI installation and UCX settings, but use different operation mappings and completion schemes. Sharing a transport framework does not isolate interface overhead: MPI RMA and two-sided operations can select different protocols and progress behavior. Similarly, comparing CUDA NCCL with SYCL oneCCL/NCCL includes compute-runtime and adapter differences.

Native and GPU-Aware MPI jobs use \texttt{srun}; oneCCL jobs use the MPI launcher matching their build. A rank wrapper derives the local rank and selects its CUDA device. Allocations request eight CPU cores per process, with binding configured to leave progress threads runnable. The oneCCL/NCCL bootstrap uses the TCP OFI provider across nodes; this does not make TCP its NCCL payload transport.

\subsection{Adapted oneCCL Paths}\label{sec:oneccl-method}

Two identifiers in Table~\ref{tab:m-evaluated-stacks} reach their backend through adapters written for this thesis, so each measures a complete adapted path rather than the backend alone. The NCCL adapter propagates oneCCL's group lifecycle to the corresponding NCCL calls, which preserves asynchronous submission: a dependent send/receive batch is issued as one group and observed through the returned event. The OSHMPI adapter maps collectives onto OpenSHMEM operations where equivalents exist and composes the remainder, building all-to-all-v from puts and synchronization, rooted reduce by discarding non-root allreduce output, and reduce-scatter from a slice of an allreduce result. Ordinary operands are staged through a grow-only symmetric arena, and because the underlying call is blocking and host-driven, pending SYCL work completes first and the returned event is already complete. That path therefore satisfies completion correctness without preserving asynchronous submission, and its staging and synchronization fall inside the timed interval because they are part of the implemented path.

The oneCCL/NVSHMEM branch is a runtime foundation rather than a backend. It establishes CUDA-backed communicator selection, UID distribution through oneCCL's key-value store, lifecycle ownership, and symmetric staging, validated on one and two nodes, while its public collective, point-to-point, barrier, and split operations remain unsupported. It contributes integration evidence for RQ2 and no performance measurement.

\subsection{Transport Availability and Verification}\label{sec:ibgda-check}

The NVSHMEM path was checked at two levels: inspection of the driver prerequisites and runtime logging of the selected transport. Table~\ref{tab:m-ibgda-check} records the relevant observations against the installation guidance \cite{nvidia_nvshmem_install_2026}. Operational identifiers and diagnostic provenance are retained in the experiment records rather than used to explain the result.

\begin{table}[htbp]
\centering
\caption{Recorded IBGDA prerequisite observations in the evaluated Leonardo environment.}
\label{tab:m-ibgda-check}
\footnotesize
\begin{tabularx}{\textwidth}{lYY}
\toprule
Prerequisite & Reference requirement & Recorded observation\\
\midrule
NVIDIA driver & At least 510.40.3 & 535.274.02\\
Peer-memory support & \texttt{nvidia\_peermem} loaded & Loaded\\
MLNX\_OFED & At least 5.0 & OFED-internal-23.04-1.1.3\\
InfiniBand interfaces & Supported \texttt{mlx5} interfaces & Four devices, one port each\\
Peer mapping override & Enabled for the documented path & Not reported in driver parameters\\
Stream memory operations & Enabled for the documented path & Disabled\\
\bottomrule
\end{tabularx}
\end{table}

A two-node diagnostic run reported the \texttt{ibrc} transport. Together with the recorded configuration, this establishes that the evaluated device-issued inter-node requests were serviced by a CPU proxy. The missing prerequisite settings are consistent with the reported lack of IBGDA support on Leonardo \cite{trotter_cpu-_2025}; an \texttt{ibrc} trace by itself would establish transport selection, not the reason another transport was unavailable. Driver requirements and alternative NIC-handler modes must be checked against the actual library release before a future GPU-posted experiment.

\subsection{Runtime Configuration and Sensitivity Controls}\label{sec:runtime-method}

MPI and OSHMPI use UCX, with HCOLL and the legacy OpenIB component disabled. Single-node transfers select shared-memory, CUDA, and self transports; multi-node transfers additionally enable reliable-connection networking, service level 1, and two rendezvous rails. The benchmark requests GPUDirect RDMA and uses \texttt{get\_zcopy} with a 1024-byte rendezvous threshold. A separate rail-count comparison and UCC on/off controls assess these choices.

UCC is enabled for standalone allreduce and the selected all-to-all configuration and disabled for \texttt{moe} and \texttt{cg\_step}. This is a fixed per-pattern policy, not a best-of-two selection at each plotted point. The UCC comparison is a separate paired configuration dataset; its results explain where the policy helps and where it incurs a penalty. Diagnostic tracing is excluded from performance aggregation because tracing and first-use connection setup change the measured interval.

NCCL retains default algorithm and protocol selection. NVSHMEM uses MPI bootstrap, the \texttt{ibrc} transport, a 1\,GiB symmetric heap, and NCCL collective dispatch in the benchmark. Its persistent inter-node neighbor kernels cap the cooperative grid at eight blocks to limit proxy request pressure. Direct OSHMPI uses a 1\,GiB heap and the UCX RMA component; oneCCL/OSHMPI uses a 2\,GiB heap, a 256\,MiB staging arena, and an 8\,MiB point-to-point slot per peer.

OSHMPI's standalone allreduce uses CUDA symmetric operands with UCC enabled. The archived \texttt{cg\_step} instead copies each device scalar to host symmetric memory and reduces it there. A separate device/UCC variant changes operand placement, allocator, and collective configuration together. These variants are kept distinct because their difference cannot be attributed solely to reduction execution.

\section{Measurement and Statistical Protocol}

\subsection{Completion Boundaries and Correctness}\label{sec:completion-method}

Timed intervals include completion rather than only submission. MPI paths include blocking communication or request completion and the required device synchronization; NCCL paths wait for stream completion; oneCCL paths wait for returned events. SHMEM paths include the completion and synchronization required by their exchange protocol: MoE's OSHMPI path, for example, closes each of dispatch and combine with a quiet, a device synchronization, and a barrier, while its NVSHMEM path closes with a quiet and one signal per peer actually exchanged with. Persistent NVSHMEM batches additionally include device-side signaling and waits. Setup, initialization, and the final correctness checks are outside the timed intervals.

An isolated halo sample completes one exchange. A steady-state sample amortizes a batch boundary over 100 exchanges while retaining the protocol's required internal synchronization. These are different execution contracts, so the amortized result is not an estimate of raw network latency. Likewise, \texttt{cg\_step}'s phase instrumentation runs separately from its unsplit timing and introduces additional synchronization.

Each implementation checks received or reduced data against its expected result and combines rank-local verdicts globally. MoE additionally validates the dispatch layout, token ownership, source-token ordering, and the deterministic payload, and then checks that combine reproduced the packed send buffer exactly. The composite benchmark checks its stencil and dot products. The archived OSHMPI scalar check validates host reduction results; it does not establish the cost of returning those results to a subsequent GPU consumer.

The adapters of Section~\ref{sec:oneccl-method} are validated at the interface level on the same principle: the criterion is that a complete batch can be submitted and its completion observed correctly, exercised through dependent send/receive batches, nested groups, deferred events, error propagation, communicator reuse, and rank agreement with the supported world team. Unsupported configurations are rejected explicitly rather than silently approximated. The halo and all-to-all benchmarks then drive the validated path under load, so interface correctness and performance retention are established by separate evidence.

\subsection{Samples, Trials, and Independent Allocations}\label{sec:samples-method}

Except for ping-pong, each rank records a local elapsed-time series, and sample $i$ is summarized using the slowest rank:
\begin{equation}
t_i=\max_{r\in\{0,\ldots,P-1\}}t_{i,r}.
\end{equation}
The arithmetic mean of this series is the trial latency. Ping-pong instead reports half of rank zero's round-trip time. Repeated iterations within a trial and repeated trials within an allocation are not independent samples of node assignment.

For the principal benchmark curves, trial means are first averaged within each allocation. The plotted center is the median of those allocation means, and shading shows their minimum--maximum range. Following the defaults of Section~\ref{sec:execution-procedure}, ping-pong, halo, allreduce, and MoE cells contain five allocations with three trials each; the retained all-to-all and \texttt{cg\_step} snapshots contain three allocations with three trials each. Ranges describe observed variation, not confidence intervals. Small differences with overlapping ranges are not treated as decisive backend advantages.

The archived phase records contain aggregated phase means rather than per-allocation phase samples. Phase tables therefore compare those means with the stored unsplit mean, while headline curves use the allocation-based median. The UCC on/off CSVs likewise retain summary means and are analyzed as a separate configuration comparison. These reporting levels are stated explicitly rather than combined into one uncertainty estimate.

\section{Communication Patterns and Application Microbenchmarks}\label{sec:benchmark-method}

\begin{table}[htbp]
\centering
\caption{Benchmark patterns and measurement ranges. Iteration counts are timed/warm-up defaults. The first four patterns are isolated operations; the last two are the application microbenchmarks.}
\label{tab:m-benchmark-patterns}
\footnotesize
\begin{tabularx}{\textwidth}{lYYl}
\toprule
Pattern & Operation & Size range & Iterations\\
\midrule
\texttt{pingpong} & Dependency-ordered peer exchange & 4\,B--16\,MiB & 100/20\\
\texttt{halo\_1d} & Periodic neighbor exchange & 16\,B--16\,MiB aggregate & 100/20\\
\texttt{allreduce} & Float32 sum & 4\,B--16\,MiB & 100/20\\
\texttt{alltoall} & Equal block per peer & 4\,B--256\,KiB per peer & 100/20\\
\midrule
\texttt{moe} & Variable-count dispatch and combine & 16\,384 tokens of 256 float32 per rank; three routing cases & 100/20\\
\texttt{cg\_step} & Halo, stencil, two FP64 sums & Global side 512--8192 & 50/10\\
\bottomrule
\end{tabularx}
\end{table}

Ping-pong uses \texttt{1n2g} and \texttt{2n1g}. Halo adds \texttt{1n4g}, \texttt{2n4g}, \texttt{4n4g}, and \texttt{8n4g}; the collectives and both application microbenchmarks also retain \texttt{1n1g} controls. Table~\ref{tab:m-pattern-coverage} gives the exact backend matrix: oneCCL/OSHMPI has standalone allreduce and all-to-all data but no ping-pong, halo, MoE, or composite-step campaign.

The topology set is deliberately not a single rank sweep. \texttt{1n2g} and \texttt{1n4g} stay on intra-node NVLink; \texttt{2n1g} isolates one InfiniBand link with no intra-node traffic; and \texttt{2n4g}, \texttt{4n4g}, and \texttt{8n4g} combine both levels at 8, 16, and 32 GPUs. Comparing \texttt{1n4g} with \texttt{2n1g} therefore separates a change of scale from a change of interconnect, which a rank-count axis alone would conflate.

Payload conventions are operation-specific. Ping-pong uses one directional payload. Halo counts two sends and two receives. Allreduce reports input bytes, with bus bandwidth normalized by $2(P-1)/P$. All-to-all reports bytes per peer on its size axis and the whole per-rank send buffer, including the self block, for aggregate bandwidth. MoE reports useful application volume, $2\,T H \cdot 4$ bytes per rank for $T$ tokens of $H$ float32 features, counting dispatch and combine for every token including those that never leave their source rank. Bandwidth is recomputed from the stated payload and latency; these different conventions are not compared as physical link rates, and the MoE figure in particular is not the traffic crossing a link.

\subsection{What the MoE Step Represents}\label{sec:moe-method}

The MoE microbenchmark measures one top-1 dispatch and its inverse combine, with no expert computation. Every rank holds $T=16{,}384$ tokens of $H=256$ float32 features. Each token is assigned to one expert rank, tokens are packed into destination-ordered blocks, and dispatch sends those variable-sized blocks to their expert ranks; combine returns each received block to its source without modifying the payload. Route generation, planning, packing, allocation, the initial barrier, and validation are outside the timed interval, so the measurement covers the redistribution itself rather than the routing decision.

Three deterministic routing cases change only the block-size distribution, not the total useful volume:

\begin{itemize}
    \item \texttt{uniform}: tokens are assigned by a uniform hash, so every peer receives a comparable block. This is the case closest to the standalone all-to-all pattern.
    \item \texttt{locality80}: approximately 80\% of tokens stay on their source rank. Off-node volume drops sharply while the number of peers does not.
    \item \texttt{hotspot80}: approximately 80\% of tokens target rank 0, producing a deliberate incast onto one GPU and its NIC.
\end{itemize}

At one rank all three cases degenerate to a local copy and are operationally equivalent; \texttt{1n1g} is therefore a packing and local-copy control, not a network measurement. Because the byte count is held fixed across cases, a latency difference between them is attributable to distribution rather than to volume. Each backend maps the exchange differently: the MPI paths issue two variable-count \texttt{MPI\_Alltoallv} calls, NCCL and oneCCL/NCCL issue grouped variable-count point-to-point exchanges, and the two SHMEM paths issue per-destination puts with local copies. The one-sided paths thus have no collective algorithm to reorganize a skewed distribution, which the \texttt{hotspot80} case is designed to expose.

MoE keeps UCC disabled. That default was inherited from the all-to-all experiment rather than established by a MoE-specific control, and MoE uses the variable-count Open MPI path, not the one the all-to-all sweep measured. The MoE results consequently describe one fixed configuration and are not a tuned MPI result.

\subsection{What the Composite Step Represents}

The composite benchmark block-distributes an $S\times S$ global grid by columns. Each iteration packs boundary columns, exchanges a nonperiodic west/east halo, computes a four-neighbor stencil and local dot products, and performs two separate global sums of one double each. The side $S$ takes values 512, 1024, 2048, 4096, and 8192. At fixed $S$, increasing rank count reduces local work; at fixed rank count, increasing $S$ increases local work quadratically and halo length linearly.

The benchmark preserves the operation sequence but omits CG vector updates, changing search directions, and convergence. Both local dot products are available before the global reductions, and the archived loop does not use the reduced scalars to form the next iterate. Thus it tests a communication/computation composition, not the complete dependency graph in Table~\ref{tab:m-cg-dependencies}. In particular, the archived staged OSHMPI path omits the host-to-device result copy required by a GPU consumer.

With phase reporting enabled, a second pass measures pack, halo, computation, and reductions separately. Synchronization at phase boundaries changes both scheduling and overhead. The difference between the phase sum and unsplit time is consequently an instrumentation-sensitive diagnostic, not a direct overlap measurement.

\section{Application Validation with aCG}\label{sec:acg-method}

\subsection{Problem and Execution Matrix}

The aCG suite provides CG and pipelined CG implementations \cite{trotter_cpu-_2025}. This campaign compares the host-initiated CG solver with MPI, NCCL, NVSHMEM, and OSHMPI communicators and separately evaluates the monolithic device-NVSHMEM solver. Both matrices use fixed METIS partitions reused across communicators: Bump\_2911 has 2,911,419 rows and 127,729,899 full symmetric nonzeros; Queen\_4147 has 4,147,110 rows and 329,499,284 full symmetric nonzeros. Runs use 1, 2, 4, 8, 16, and 32 GPUs.

The manufactured solution uses seed 101, relative residual tolerance $10^{-6}$, zero absolute tolerance, and a limit of 100,000 iterations. The retained archive contains 507 completed solver logs: 354 host-communicator trials, 90 device-NVSHMEM trials, and 63 single-GPU controls. These counts describe available completed records; incomplete attempts do not become successful trials by omission. Backup copies are excluded from the analysis.

Two properties of this matrix constrain what the cross-level comparison can claim. First, the measured OSHMPI communicator implements the halo only: it allocates CUDA symmetric buffers, exchanges remote offsets at setup, and completes its blocking puts inside halo-begin, before local SpMV, while its scalar reductions use CUDA-aware \texttt{MPI\_Allreduce}. The application's OSHMPI label therefore identifies a halo implementation, not an OpenSHMEM scalar-reduction measurement. Second, the solver departs from the benchmark defaults in a 16,384-byte UCX rendezvous threshold and in disabling NCCL dispatch beneath NVSHMEM; separate sensitivity controls assess both, so the archived default NVSHMEM reduction is not treated as a configuration-matched prediction. The monolithic device solver additionally reorganizes computation and synchronization, and is evaluated as an execution-model variant rather than a communicator substitution.

\subsection{Solver Metrics and Prediction Accuracy}

Solver time is the time-to-solution measure. Time per iteration and analytical throughput help distinguish execution cost from changes in convergence. Following the reference study, the primary scaling view selects the fastest residual-valid completed trial per matrix, backend, and rank count. For standard CG, the analytical work is
\[
W(K)=2(K+1)\,\mathrm{nnz}+2(5K+1)n,
\]
where $K$ is the executed iteration count, $n$ the row count, and $\mathrm{nnz}$ the full nonzero count. Per-GPU throughput is $W(K)/(P T)$; it is not a hardware-counter measurement. All retained trial times remain available, and MPI variability is shown by allocation medians and ranges rather than a line connecting rare fast trials.

Correctness is checked from the recorded residual relative to the initial residual and the requested tolerance, together with successful completion. Manufactured-solution error is retained as a separate numerical diagnostic. Missing correctness fields are distinguished from failures. The parser does not infer convergence solely from a process exit code.

Phase timers have backend-dependent meanings. The halo timer brackets the wait remaining at halo-end, whereas OSHMPI's blocking communication appears mainly in \emph{other}. NVSHMEM may include communication in its SpMV interval, and allreduce timing includes arrival skew. These columns cannot simply be added into an overlap-free network-cost model.

For the scalar cross-level comparison, half of the stored two-reduction \texttt{cg\_step} phase at side 512 is compared with aCG's accumulated allreduce time divided by its recorded operation count. The latter includes operations outside the steady loop: a run with $K$ iterations records $2K+2$ reductions. MPI is paired with the MPI scalar calls in the OSHMPI-halo solver, and NCCL with the NCCL solver. Mean absolute percentage discrepancy over the five matched rank counts is
\[
\mathrm{MAPE}=\frac{100}{5}\sum_{P\in\{2,4,8,16,32\}}
\left|\frac{\widehat{t}_P}{t_P}-1\right|.
\]
This is a descriptive phase-cost check across increasing scale, not a fitted or held-out prediction of total solver time. Matching eight bytes in the float32 microbenchmark alone does not establish equivalence to a one-double sum; the composite-to-solver pairing matches datatype as well as payload.

The CUDA--SYCL question is answered by a separate report rather than by this campaign. It compares solver-time medians on Bump\_2911 at 1, 4, 8, and 16 GPUs over three repeats, with MPI supplying distributed execution, and uses medians on both sides rather than the fastest-trial convention above. Its CUDA parser does not populate the convergence flag, so its zero converged-count field represents missing classification rather than demonstrated nonconvergence. Until those logs are reprocessed, the ratios are retained as supporting implementation timings, not a verified equal-accuracy comparison, and they do not separate kernel, memory, and orchestration contributions.

\section{Reproducibility and Evidence Boundaries}\label{sec:reproducibility}

The reconciled archive in \path{data/main-campaign/} contains benchmark allocation summaries, all retained aCG trial records, selected solver results, and reduction-prediction pairs. A manifest records source-file hashes and the repository revisions inspected during import. Those revisions identify the imported archive, not necessarily the source used for each measured run. \path{tools/rebuild_main_evidence.py} regenerates the data and complete table blocks; \path{tools/plot_main_evidence.py} renders the figures. The accompanying README specifies the aggregation and exclusion policies.

The MoE campaign is a later benchmark snapshot than the reconciled inputs, so it is imported separately by \path{tools/import_moe_evidence.py} into its own summary, allocation, and manifest files rather than merged into the shared benchmark summary. One archive therefore corresponds to one snapshot, and the shared aggregation policy is reused rather than reimplemented. The archive is explicitly a mixed snapshot for the same reason. The allreduce sweep was re-run and re-exported after the original import, and the reported allreduce values are that later campaign; the \texttt{cg\_step} sweep was re-run only at global side 512, so the wider archived sweep is retained rather than replaced. The import records which patterns were retained and carries their original source hashes forward, and it refuses, before writing anything, an input that would narrow the archive. The exported CSV files, not the current source tree, are the evidence for the retained patterns.

Auxiliary rail, scalar-placement, forced-algorithm, and runtime-dispatch controls survive partly as analysis records rather than complete per-trial datasets. Their reported effects are identified as such and are not assigned reconstructed error bars. The original IBGDA job and node identifiers are retained in \path{data/experiment-notes.md}. Preliminary NCCL transcription and plotting provenance are recorded separately in \path{data/nccl-preliminary.md}.

The next measurements are defined in \path{experiments/targeted-controls.md}. Their priority relative to matched initiation, batching, and configuration controls is recorded in \path{experiments/missing-experiments-report.md}. These documents describe prospective work and do not turn current unmeasured source changes into experimental evidence.
